% Standalone LaTeX section for the paper.
% This snippet is written to drop into the original manuscript after the
% execution-model section, replacing the earlier evaluation-plan placeholder.

\section{Evaluation}
\label{sec:evaluation}

This section reports empirical results for the Tiny SDD LLM benchmark after
instantiating the Spec-Driven Development workflow described in
Section~\ref{sec:exec}. The evaluation separates \emph{easy} and \emph{hard}
scenario families because the benchmark now contains two qualitatively different
task distributions: simpler, mostly single-hop repository edits where standard
test-gated generation is already near ceiling, and harder multi-file tasks with
cross-module dependencies and stricter governance requirements.

\subsection{Experimental setup}

The benchmark contains $12$ scenario packages under
\texttt{dataset\_sdd/}. We group them as follows.

\begin{itemize}[leftmargin=*]
\item \textbf{Easy scenarios:} \texttt{scenario2--scenario9} and
\texttt{scenario12}, for a total of $9$ scenarios and $108$ task runs per mode.
These scenarios contain $12$ tasks each and represent relatively direct feature
additions and bug fixes with limited cross-file coordination.
\item \textbf{Hard scenarios:} \texttt{scenario1}, \texttt{scenario10}, and
\texttt{scenario11}, for a total of $3$ scenarios and $45$ task runs per mode.
These scenarios contain $15$ tasks each and introduce richer app structure,
additional source files, stricter Architecture Cards, and more frequent
cross-file constraints.
\end{itemize}

Each task was executed once per mode in the current reported runs
($1$ trial per task). The evaluated modes are:
\texttt{tests\_only}, \texttt{baseline}, and \texttt{full\_sdd} on all
scenarios, plus \texttt{full\_sdd\_no\_operator} on the hard split as an
ablation. The ablation keeps the SDD verifier active but disables the
deterministic compliance operator that injects required reads, trace witnesses,
and test touches before verification. This isolates whether the benefit comes
from the gate alone or from the full governed repair loop.

\paragraph{Metrics.}
We report:
\textit{success} (accepted runs / total runs),
\textit{compliant success} (accepted runs with zero remaining policy violations),
\textit{average policy violations},
\textit{average rejections},
\textit{average retries},
\textit{unresolved rejections} (share of runs rejected at least once and still
ending in failure), and
\textit{agent error rate} (share of runs ending in upstream agent/runtime error).
In this prototype, policy violations are lightweight post-run governance checks
over the final candidate, including issues such as PII leakage, missing required
reads before sensitive edits, missing test updates for public API changes, and
other Architecture Card violations.

\subsection{Results on easy scenarios}

Table~\ref{tab:easy-results} shows that the easy split is close to ceiling for
all ungated methods. Both \texttt{tests\_only} and \texttt{baseline} achieve
$1.00$ functional success. Their compliance is also already high
($0.98$ and $0.97$, respectively), which means these tasks do not strongly
stress governance failures. Under this easier workload, \texttt{full\_sdd}
slightly reduces functional success to $0.95$ and introduces modest repair-loop
activity ($0.28$ average rejections, $0.23$ average retries). This behavior is
consistent with a stricter gate operating in a low-headroom regime: when the
task distribution is simple, there is less room for compliance gains and more
chance that a conservative repair loop trades a small amount of throughput for
assurance.

\begin{table}[t]
\centering
\small
\begin{tabular}{lccccc}
\toprule
\textbf{Mode} & \textbf{Success} & \textbf{Compliant} & \textbf{Avg.\ Viol.} & \textbf{Avg.\ Rej.} & \textbf{Avg.\ Retry} \\
\midrule
\texttt{tests\_only} & 1.00 & 0.98 & 0.02 & 0.00 & 0.00 \\
\texttt{baseline} & 1.00 & 0.97 & 0.03 & 0.00 & 0.00 \\
\texttt{full\_sdd} & 0.95 & 0.95 & 0.05 & 0.28 & 0.23 \\
\bottomrule
\end{tabular}
\caption{Results on the easy split (\mbox{$9$ scenarios, $108$ task runs per mode}).}
\label{tab:easy-results}
\end{table}

\subsection{Results on hard scenarios}

The hard split is where SDD becomes materially useful. Table~\ref{tab:hard-results}
shows a clear separation between functional success and compliant success for the
non-SDD baselines. \texttt{baseline} reaches $0.93$ functional success, but only
$0.58$ compliant success; in absolute terms, it accepts $42/45$ tasks, yet only
$26/45$ are both accepted and governance-compliant. \texttt{tests\_only} is
weaker still, with $40/45$ accepted but only $22/45$ compliant.

By contrast, \texttt{full\_sdd} preserves the same functional success as
\texttt{baseline} ($0.93$, or $42/45$ accepted) while raising compliant success
to $0.93$ ($42/45$ compliant). Relative to \texttt{baseline}, this is a
$35$-point improvement in compliant success and an $86\%$ reduction in average
policy violations (from $0.47$ to $0.07$). The additional enforcement overhead
is limited: \texttt{full\_sdd} averages only $0.22$ rejections and
$0.16$ retries per run, with unresolved rejections in $0.07$ of runs.

\begin{table}[t]
\centering
\small
\begin{tabular}{lccccc}
\toprule
\textbf{Mode} & \textbf{Success} & \textbf{Compliant} & \textbf{Avg.\ Viol.} & \textbf{Avg.\ Rej.} & \textbf{Avg.\ Retry} \\
\midrule
\texttt{tests\_only} & 0.89 & 0.49 & 0.58 & 0.00 & 0.00 \\
\texttt{baseline} & 0.93 & 0.58 & 0.47 & 0.00 & 0.24 \\
\texttt{full\_sdd\_no\_operator} & 0.27 & 0.27 & 1.36 & 2.24 & 1.51 \\
\texttt{full\_sdd} & 0.93 & 0.93 & 0.07 & 0.22 & 0.16 \\
\bottomrule
\end{tabular}
\caption{Results on the hard split (\mbox{$3$ scenarios, $45$ task runs per mode}).}
\label{tab:hard-results}
\end{table}

\subsection{Ablation: gate-only SDD is not enough}

The hard-split ablation is the most informative result in the study.
\texttt{full\_sdd\_no\_operator} performs poorly on every major metric:
$0.27$ success, $1.36$ average policy violations, $2.24$ average rejections,
and $0.73$ unresolved rejections. In absolute terms, only $12/45$ hard tasks
are both accepted and compliant. This result is important because it shows that
the value of SDD does not come from a strict gate alone. A gate without
operator-assisted repair is too brittle: it rejects unsafe proposals, but often
fails to steer the agent back toward an acceptable patch.

The empirical implication is that SDD should be understood as a
\emph{governed generation-and-repair workflow}, not merely as a post-hoc
verification barrier. The Architecture Card, trace checks, and capability bounds
need to be coupled with targeted repair guidance and deterministic compliance
assists if the system is expected to remain usable on harder repository tasks.

\subsection{Discussion}

Taken together, the easy and hard splits suggest two different operating regimes.
On easy tasks, repository edits are simple enough that standard test-gated
generation already performs near ceiling, so SDD has limited room to improve
compliance and may incur a small throughput penalty. On hard tasks, however,
there is a substantial \emph{assurance gap}: many outputs that pass tests remain
policy-unsafe. In that regime, \texttt{full\_sdd} closes the gap without
materially sacrificing task completion.

This pattern is aligned with the motivating claim of the paper. The main value
of SDD is not that it universally raises functional success on all tasks. Its
value is that, when repository changes become multi-step and governance-sensitive,
it can keep functional performance competitive while sharply improving the rate
of accepted \emph{and compliant} outcomes.

\subsection{Threats to validity}

These results should be interpreted with several limitations in mind.
First, the current reported runs use a single trial per task, so they do not yet
support confidence intervals over model stochasticity. Second, the benchmark is
synthetic and compact; while it captures cross-file dependencies and governance
constraints, it does not match the scale or organizational complexity of a
production Adobe repository. Third, the easy split was evaluated before the
\texttt{full\_sdd\_no\_operator} ablation was introduced, so the ablation results
currently apply only to the hard scenarios. Fourth, the \textit{policy violations}
metric is a lightweight proxy for governance non-compliance rather than a formal
proof obligation. These limitations do not invalidate the observed hard-split
effect, but they do bound the strength of any external claims.

\paragraph{Summary of findings.}
The current prototype supports three concrete conclusions:
\begin{enumerate}[leftmargin=*]
\item On easy tasks, all methods are already near ceiling, and SDD introduces a
small but measurable overhead.
\item On hard tasks, \texttt{full\_sdd} matches the best functional baseline
($0.93$ success) while substantially improving compliant success
($0.93$ vs.\ $0.58$ for \texttt{baseline}).
\item The operator-assisted repair loop is essential; gate-only SDD is too
brittle to be practically competitive.
\end{enumerate}
